% =====================================================================
%  estrategia_market_making.tex
%  Documentación técnica — Bot de Market Making adaptativo
%  Instrumento: XRPUSDC — Binance Futures (USDS-M)
%  Compilar: pdflatex -interaction=nonstopmode -output-directory=doc
%            doc/estrategia_market_making.tex  (ejecutar 2 veces)
% =====================================================================
\documentclass[11pt,a4paper]{article}

\usepackage[utf8]{inputenc}
\usepackage[T1]{fontenc}
\usepackage[spanish,es-noshorthands]{babel}
\usepackage{amsmath,amssymb}
\usepackage[margin=2.4cm]{geometry}
\usepackage{booktabs}
\usepackage{longtable}
\usepackage{array}
\usepackage[hidelinks]{hyperref}

\setlength{\parskip}{0.4em}
\setlength{\parindent}{0pt}

\newcommand{\codigo}[1]{\texttt{#1}}
\newcommand{\symbolname}[1]{\textsc{#1}}

\title{\bfseries Bot de \emph{Market Making} Adaptativo\\[0.5em]
{\large XRPUSDC — Binance Futures (USDS-M)}\\[1.2em]
{\normalsize Documentación técnica de la estrategia v1.0}}
\author{Proyecto \emph{alpha-driven}}
\date{10 de agosto de 2026}

\begin{document}

\maketitle

% =====================================================================
% RESUMEN EJECUTIVO
% =====================================================================
\section*{Resumen ejecutivo}
\addcontentsline{toc}{section}{Resumen ejecutivo}

Este documento describe la estrategia de \emph{market making} adaptativo
implementada en el repositorio \codigo{alpha-driven}: un bot que cotiza
simultáneamente un precio de compra (bid) y uno de venta (ask) alrededor
del precio medio del libro de órdenes del par \codigo{XRPUSDC} en
Binance Futures (USDS-M), con tamaño y \emph{spread} ajustados por la
volatilidad observada y el inventario acumulado. El objetivo es capturar
el \emph{spread} como ingreso neto de comisiones, \emph{funding} y
deslizamiento (NetPnL $> 0$, véase §6).

La implementación cumple las reglas no negociables de la Sección 0 de la
especificación (§1): testnet primero, presupuestos de riesgo definidos
como constantes, kill switch probado explícitamente y detención ante
situaciones irreconciliables. La decisión de diseño aprobada habilita la
corrida integrada \emph{dry-run} en Nivel de exposición 0 sin enviar
órdenes reales. A la fecha de este documento la estrategia operó su
primera corrida en mainnet (autorización humana explícita, §0.2) con
resultados de reconciliación exitosos; los presupuestos de riesgo siguen
marcados como propuestas pendientes de confirmación (§0.4).

\subsection*{Tablero de estado (al 10/08/2026)}
\begin{center}
\begin{tabular}{ll}
\toprule
\textbf{Componente} & \textbf{Estado} \\
\midrule
Risk Engine (tests unitarios) & 13/13 OK \\
Execution Engine (tests unitarios) & 14/14 OK \\
Alpha Model (tests unitarios) & 18/18 OK \\
Kill switch (prueba explícita) & 4/4 OK \\
Primera corrida mainnet & Operada: 8 órdenes reales, 4.9 XRP, exit 0 \\
Presupuestos de riesgo & Propuestos, pendientes de confirmación (§0.4) \\
\bottomrule
\end{tabular}
\end{center}

\tableofcontents
\newpage

% =====================================================================
% 1. OBJETIVO Y REGLAS NO NEGOCIABLES
% =====================================================================
\section{Objetivo y reglas no negociables (§0)}

\subsection{Objetivo}
El sistema debe generar ganancia neta consistente operando como
\emph{market maker} en el par \codigo{XRPUSDC} de Binance Futures,
adaptando su comportamiento a la volatilidad del mercado y al inventario
acumulado, sin superar nunca los límites de riesgo definidos en
\codigo{config.py}.

\subsection{Reglas no negociables (§0)}
Las siguientes reglas tienen prioridad sobre el resto de la
especificación y sobre cualquier decisión operativa:

\begin{enumerate}
  \item \textbf{Testnet primero.} Mainnet SOLO con autorización humana
  explícita. No se llama \codigo{init\_client(real=True)} ni
  \codigo{set\_testnet(False)} sin esa autorización (§0.1).
  \item \textbf{Detenerse y preguntar} ante: primera orden en mainnet,
  subir el Nivel de exposición, kill switch activado, falta de una
  función de API (no escribir implementación paralela), posición no
  reconciliable o credenciales nuevas (§0.2).
  \item \textbf{Presupuesto de riesgo} como constantes en
  \codigo{config.py}, nunca números improvisados (§0.4):
  \begin{equation*}
    \codigo{MAX\_DAILY\_LOSS\_USDC} \qquad
    \codigo{MAX\_POSITION\_NOTIONAL\_USDC} \qquad
    \codigo{MAX\_DRAWDOWN\_PCT} \qquad
    \codigo{MAX\_LEVERAGE\_USED}
  \end{equation*}
  Si no están definidos: proponer valores conservadores, marcarlos
  ``pendiente de confirmación'' y no operar con fondos reales.
  \item \textbf{Credenciales} (§0.5): nunca loguearlas, imprimirlas,
  commitearlas ni rotarlas sin preguntar. \codigo{.env} está en
  \codigo{.gitignore}.
  \item \textbf{Commits chicos y frecuentes} (§0.3) y \codigo{STATUS.md}
  actualizado cada 15--20 minutos.
\end{enumerate}

\subsection{Decisión de diseño aprobada: Nivel 0 / dry-run}
Cuando $\codigo{EXPOSURE\_LEVEL}=0$ el bot NO queda bloqueado en tamaño
cero: puede simular órdenes usando la constante de simulación
$\codigo{SIMULATION\_QUOTE\_MULTIPLIER}$ de \codigo{config.py} en lugar
del multiplicador $0.0$ del nivel 0. Esto habilita la corrida integrada
\emph{dry-run}. Los niveles reales ($\geq 1$) NO se modifican y subirlos
sigue requiriendo autorización humana (§0.2). En Nivel 0 jamás se envían
órdenes reales (§0.1/§21).

% =====================================================================
% 2. ARQUITECTURA
% =====================================================================
\section{Arquitectura}

\subsection{Interfaz única con Binance}
\codigo{API\_binance\_futuros.py} es la única interfaz con Binance
Futures (§1). No se duplican funciones que ya existen; ante una función
faltante se debe detener y preguntar (§0.2). Sus características:

\begin{itemize}
  \item Inicialización: \codigo{init\_client(real: bool = False)}.
  Testnet usa el modo nativo de \codigo{python-binance}; mainnet fuerza
  URLs reales.
  \item Los métodos devuelven \codigo{BinanceAPIException} como
  resultado en vez de lanzarla: el llamador debe chequear el tipo de
  retorno.
  \item \codigo{\_call\_with\_retry(func, retries=2, delay=0.5, ...)}
  envuelve las llamadas críticas.
  \item Sincronización de tiempo: \codigo{adjust\_client\_time()} más un
  hilo daemon cada 300~s que fija \codigo{client.timestamp\_offset}.
\end{itemize}

\subsection{WebSocket para datos en tiempo real}
El flujo de datos es por WebSocket; REST se usa solo para configuración,
estado, órdenes, recuperación y reconciliación (§2). URLs: real
\codigo{wss://fstream.binance.com/public/ws}, testnet
\codigo{wss://stream.binancefuture.com/public/ws}. Reintento cada 5~s
con \codigo{ping\_interval=20}. Si se superan los reintentos, se activa
el kill switch (§13).

Los módulos expuestos y sus callbacks:
\begin{itemize}
  \item \codigo{websocket\_bookticker.py} $\rightarrow$
  \codigo{BookTickerWebSocket(symbol, on\_bookticker, real=False)},
  stream $\{\text{symbol}\}$@bookTicker. Callback:
  \codigo{on\_bookticker(bid\_price, bid\_qty, ask\_price, ask\_qty, timestamp)}.
  \item \codigo{websocket\_depth.py} $\rightarrow$
  \codigo{DepthWebSocket(symbol, on\_depth, real=False, depth\_levels=5)},
  stream $\{\text{symbol}\}$@depth\{levels\}@100ms. Callback:
  \codigo{on\_depth(bids, asks, depth, timestamp)} con \codigo{bids} y
  \codigo{asks} como listas $[\text{precio},\,\text{qty}]$.
  \item \codigo{websocket\_trades.py} $\rightarrow$
  \codigo{TradeWebSocket(symbol, on\_trade, real=False)}, stream
  $\{\text{symbol}\}$@trade. Callback:
  \codigo{on\_trade(price, qty, is\_buyer\_maker, timestamp)}.
\end{itemize}

\subsection{Flujo de datos}
El estado de mercado (\codigo{market\_state.py}) consume los tres
streams, mantiene libro, últimos precios y régimen de volatilidad con
\codigo{threading.Lock}. El \codigo{market\_maker.py} es el orquestador:
cada ciclo calcula cotizaciones, consulta el Risk Engine y delega el
envío al \codigo{execution\_engine.py}. El
\codigo{inventory\_manager.py} lleva el inventario y el PnL realizado.

\begin{center}
\begin{tabular}{ll}
\toprule
\textbf{Capas} & \textbf{Módulos} \\
\midrule
Datos & \codigo{websocket\_bookticker/depth/trades.py}, \codigo{market\_state.py} \\
Lógica & \codigo{alpha\_model.py} (cotizaciones), \codigo{inventory\_manager.py} \\
Riesgo & \codigo{risk\_engine.py} (rechazos y kill switch) \\
Ejecución & \codigo{execution\_engine.py} \\
Orquestación & \codigo{market\_maker.py}, \codigo{run\_dry\_run.py}, \codigo{run\_mainnet.py} \\
Configuración & \codigo{config.py} \\
\bottomrule
\end{tabular}
\end{center}

\subsection{Instrumento}
El instrumento es \codigo{XRPUSDC}, con el símbolo como parámetro de
configuración (no hardcodeado). Las specs reales del símbolo
(precisiones, tick/step size, min quantity/notional) se consultan
dinámicamente — nunca se asumen (§3).

% =====================================================================
% 3. ESTRATEGIA DE COTIZACIÓN
% =====================================================================
\section{Estrategia de cotización}

\subsection{Régimen de volatilidad}
La volatilidad se estima como el desvío estándar muestral de los
retornos logarítmicos del precio medio del libro:

\begin{equation}
  r_t = \ln\!\left(\frac{P_t^{\mathrm{mid}}}{P_{t-1}^{\mathrm{mid}}}\right),
  \qquad
  \hat{\sigma} = \sqrt{\frac{1}{n-1}\sum_{t=1}^{n}\left(r_t - \bar{r}\right)^2}
\end{equation}

donde $n$ es el número de retornos en la ventana
(\codigo{VOLATILITY\_WINDOW\_SEC}) y cada retorno corresponde al
intervalo de muestreo \codigo{SAMPLING\_INTERVAL\_SEC} (5~s). La
volatilidad NO se anualiza (§0.6). Como la ventana puede tener huecos,
el régimen se reescala a un intervalo de referencia
\codigo{INTERVALS\_PER\_MINUTE} $= 12$:

\begin{equation}
  \hat{\sigma}_{\min} = \hat{\sigma}\cdot\sqrt{12},
  \qquad
  \hat{\sigma}_{\mathrm{ref}}
  = \hat{\sigma}\cdot\sqrt{\frac{\codigo{SAMPLING\_INTERVAL\_SEC}}
  {\text{intervalo medio observado}}}
\end{equation}

\emph{Convención crítica (§0.6):} cada factor $\sqrt{T}$ se aplica
exactamente una vez. La volatilidad \emph{efectiva} que escala las
distancias de cotización es:

\begin{equation}
  \sigma_{\mathrm{efectiva}}
  = \hat{\sigma}_{\min}\cdot\sqrt{\codigo{CYCLE\_INTERVAL\_SEC}}
  = \hat{\sigma}\cdot\sqrt{12}\cdot\sqrt{\codigo{CYCLE\_INTERVAL\_SEC}}
\end{equation}

Aplicar $\sqrt{T}$ dos veces (volatilidad anualizada y luego reescalada
por el ciclo) es un error conocido del proyecto (§0.6) y está cubierto
por tests del alpha model.

\subsection{Precio medio y spread objetivo}
El precio medio del libro (mid) se calcula con el mejor bid y el mejor
ask:

\begin{equation}
  P_{\mathrm{mid}} = \frac{P_{\mathrm{bid}} + P_{\mathrm{ask}}}{2}
\end{equation}

El \emph{spread} objetivo es proporcional a la volatilidad efectiva:

\begin{equation}
  S = k\cdot\sigma_{\mathrm{efectiva}}
\end{equation}

donde $k$ es un factor de escala de \codigo{config.py}. A mayor
volatilidad, mayor spread para compensar el riesgo de inventario.

\subsection{Precios de cotización (bid/ask) con skew}
El modelo ajusta los precios alrededor del mid con un sesgo
(\emph{skew}) dependiente del inventario: si se tiene inventario largo
(comprado), la cotización de venta se acerca al mid (incentivo a
desacumular) y la de compra se aleja; simétricamente para inventario
corto:

\begin{equation}
  P_{\mathrm{bid}} = P_{\mathrm{mid}} - \frac{S}{2} + \delta_{\mathrm{inv}},
  \qquad
  P_{\mathrm{ask}} = P_{\mathrm{mid}} + \frac{S}{2} + \delta_{\mathrm{inv}}
\end{equation}

con el sesgo $\delta_{\mathrm{inv}}$ calculado por
\codigo{alpha\_model.py} en función del inventario normalizado
(convención de signo: inventario $> 0$ es LARGO, $< 0$ es CORTO; en el
libro, \codigo{BUY} suma cantidad e \codigo{SELL} resta). Los precios se
redondean al \emph{tick size} del símbolo consultado dinámicamente (§3).

\subsection{Tamaños de orden y clamp}
El tamaño base de cada lado es \codigo{BASE\_ORDER\_SIZE\_XRP}. Cuando
ya existe inventario, el tamaño del lado que reduce inventario
(\emph{reduce}) se calcula como:

\begin{equation}
  q_{\mathrm{reduce}}
  = \min\!\bigl(\codigo{MAX\_ORDER\_SIZE\_XRP},\;
    |\text{inventario}| + \codigo{BASE\_ORDER\_SIZE\_XRP}\bigr)
\end{equation}

El clamp contra \codigo{MAX\_ORDER\_SIZE\_XRP} es necesario porque una
orden de reducción no debe superar el tamaño máximo de orden del
símbolo: durante la primera corrida mainnet se observó el error de
Binance \codigo{order\_size\_exceeded} con inventario de 4.9~XRP, que
bloqueaba la reducción de posición. Con el clamp, el lado reduce queda
acotado a \codigo{MAX\_ORDER\_SIZE\_XRP} y la posición se puede
desacumular sin error (fix 632f9c3).

El tamaño total expuesto por lado queda limitado por el Risk Engine
(§4); en Nivel 0 el multiplicador de cotización es
$\codigo{SIMULATION\_QUOTE\_MULTIPLIER}$ (decisión de diseño aprobada,
§1.3).

% =====================================================================
% 4. RISK ENGINE
% =====================================================================
\section{Risk Engine}

\subsection{Rol}
El Risk Engine (\codigo{risk\_engine.py}) es un módulo independiente que
lee los límites de \codigo{config.py} y decide si una operación es
admisible. Toda orden propuesta pasa por él antes de enviarse; ante un
rechazo (\codigo{allowed=False}) la orden no se envía. Debe ser probado
con tests unitarios que verifiquen los rechazos (§12) — actualmente
13/13 OK.

\subsection{Presupuestos como desigualdades}
Los límites son constantes en \codigo{config.py} (§0.4), propuestas
conservadoras pendientes de confirmación humana:

\begin{center}
\begin{tabular}{ll}
\toprule
\textbf{Límite} & \textbf{Desigualdad} \\
\midrule
Pérdida diaria & $\text{pérdida del día} \leq \codigo{MAX\_DAILY\_LOSS\_USDC}$ \\
Nocional de posición & $|q\cdot P| \leq \codigo{MAX\_POSITION\_NOTIONAL\_USDC}$ \\
Drawdown & $\text{drawdown} \leq \codigo{MAX\_DRAWDOWN\_PCT}$ \\
Apalancamiento & $\text{apalancamiento} \leq \codigo{MAX\_LEVERAGE\_USED}$ \\
Tamaño de orden & $q \leq \codigo{MAX\_ORDER\_SIZE\_XRP}$ \\
\bottomrule
\end{tabular}
\end{center}

Si no estuvieran definidos, el bot no operaría con fondos reales.

\subsection{Niveles de exposición}
El multiplicador de tamaño se controla con
$\codigo{EXPOSURE\_LEVEL} \in \{0,1,2,3\}$:

\begin{center}
\begin{longtable}{cll}
\toprule
\textbf{Nivel} & \textbf{Modo} & \textbf{Comportamiento} \\
\midrule
\endhead
$0$ & Dry-run & Simulación en memoria con
$\codigo{SIMULATION\_QUOTE\_MULTIPLIER}$. Jamás envía órdenes reales. \\
$1$ & Testnet real & Tamaño reducido; operación real solo en testnet. \\
$2$ & Testnet completo & Tamaño nominal con los límites de config. \\
$3$ & Mainnet & Operación real con fondos. Exige autorización humana explícita (§0.1/§0.2). \\
\bottomrule
\end{longtable}
\end{center}

Subir de nivel de exposición real requiere autorización humana (§0.2).

\subsection{Kill switch (§13)}
El kill switch se activa ante cualquiera de estas condiciones:
pérdida excesiva, comportamiento anómalo, WebSocket desconectado, estado
inconsistente, órdenes desconocidas, posición no reconciliada, error de
API repetitivo, precio anómalo, pérdida de sincronización o exposición
inesperada.

Procedimiento de emergencia:
\begin{enumerate}
  \item Cancelar TODAS las órdenes abiertas
  (\emph{cancel all open orders}).
  \item Reconciliar la posición real contra el estado interno.
  \item Nunca asumir que cancelar órdenes implica no tener posición:
  el inventario se verifica contra la cuenta.
\end{enumerate}

El kill switch debe probarse explícitamente antes de confiar en él
(§13) — cubierto por 4 tests OK. En Nivel 0, su activación es análoga
pero no toca la cuenta real.

% =====================================================================
% 5. EJECUCIÓN Y RECONCILIACIÓN
% =====================================================================
\section{Ejecución y reconciliación}

\subsection{Órdenes maker}
Todas las órdenes son \emph{maker} (post-only) para capturar el spread
sin pagar comisiones de taker:

\begin{itemize}
  \item \codigo{timeInForce='GTX'} (post-only) en cada orden.
  \item \codigo{positionSide='BOTH'} (modo one-way).
  \item \codigo{reduceOnly=True} cuando corresponde (lado que reduce
  inventario).
  \item Tamaño y precio validados contra \emph{min quantity}, \emph{min
  notional}, \emph{tick size} y \codigo{PERCENT\_PRICE} del símbolo.
\end{itemize}

La verificación de que el precio cotizado es realmente \emph{maker}
(no cruza el libro, §10) se hace antes de cada envío.

\subsection{Identificación idempotente de órdenes}
Cada orden lleva un \codigo{client\_order\_id} único por lado:

\begin{verbatim}
MM-<lado>-<timestamp>-<contador>   # ej.: MM-BUY-1723300000123-1
\end{verbatim}

Los IDs de bid y ask son independientes y se purgan al confirmar la
muerte de la orden (§0.6: compartir IDs bid/ask sin purgar genera
``órdenes fantasma'' acumuladas). El contador garantiza idempotencia
ante reinicios: al reenviar, el mismo ID no duplica la orden.

\subsection{Detección de fills y reconciliación}
La detección de \emph{fills} ocurre en el \emph{refresh} de órdenes
abiertas: si una orden desaparece del libro (vía
\codigo{sync\_removed}) y la posición de la cuenta creció en el mismo
sentido, se registra un \emph{fill} real y se actualiza el inventario
(en modo \emph{dry-run}, el fill es simulado en memoria).

La reconciliación distingue dos casos que antes se confundían:

\begin{center}
\begin{tabular}{ll}
\toprule
\textbf{Caso} & \textbf{Tratamiento} \\
\midrule
Posición vacía (\codigo{0}) & Posición legítima; pre-flight continúa \\
Error de API (\codigo{None}) & Falla de consulta; detenerse y reintentar \\
\bottomrule
\end{tabular}
\end{center}

Este fix (345c7a7) destrabó el pre-flight de mainnet con cuenta limpia.
Ante una posición que no se puede reconciliar, la regla §0.2 ordena
detenerse y preguntar.

\subsection{Concurrencia y rate-limit}
El \codigo{execution\_engine.py} protege el estado compartido con un
\codigo{threading.Lock} y no sincroniza estado/inventario en cada ciclo
sin control de frecuencia (rate-limit, §0.6). Los envíos críticos pasan
por \codigo{\_call\_with\_retry}.

% =====================================================================
% 6. MÉTRICAS DE RENDIMIENTO
% =====================================================================
\section{Métricas de rendimiento}

\subsection{NetPnL}
El criterio de éxito es ganancia neta:

\begin{equation}
  \text{NetPnL} = \text{GrossPnL} - \text{fees} - \text{funding} - \text{slippage}
\end{equation}

($\S$18 de la especificación). Los componentes:
\begin{itemize}
  \item \textbf{GrossPnL}: diferencia de precios de las operaciones
  cerradas, con PnL realizado por media móvil en
  \codigo{inventory\_manager.py}.
  \item \textbf{Fees}: comisiones de \emph{maker}/\emph{taker}
  acumuladas por separado (\codigo{total\_fees}) y neteadas aquí.
  \item \textbf{Funding}: tasa por 8~h tomada de
  \codigo{FUNDING\_RATE\_PER\_8H} de \codigo{config.py} (el bot no
  consulta \codigo{get\_funding\_rate}).
  \item \textbf{Slippage}: deslizamiento estimado entre el precio
  cotizado y el precio de ejecución.
\end{itemize}

\subsection{Una operación ganadora NO es evidencia (§19)}
La rentabilidad se evalúa con \codigo{NetPnL > 0} acumulado, no con
operaciones individuales. Una operación ganadora puede ser ruido o
compensar pérdidas previas; solo el PnL neto agregado valida la
estrategia.

% =====================================================================
% 7. OPERACIONES
% =====================================================================
\section{Operaciones}

\subsection{Watchdog}
El bot corre un watchdog que vigila la actividad del ciclo de
cotización. El timeout depende del número de ciclos configurado
(\codigo{--cycles}) con un piso:

\begin{equation}
  W_{\mathrm{watchdog}}
  = \max\!\bigl(125~\text{s},\; 6\cdot\codigo{cycles} + 30~\text{s}\bigr)
\end{equation}

Si el ciclo no reporta actividad dentro de $W_{\mathrm{watchdog}}$, el
proceso termina con error (y, en niveles reales, se activa el kill
switch si corresponde).

\subsection{Pre-flight en mainnet}
Antes de operar, \codigo{run\_mainnet.py} valida condiciones y
credenciales y tolera una posición pre-existente como \emph{round trip}
si se mantiene dentro de los límites:

\begin{equation}
  |q_{\mathrm{pos}}\cdot P_{\mathrm{mark}}|
  \leq \codigo{MAX\_POSITION\_NOTIONAL\_USDC}
  \quad\text{y}\quad
  |q_{\mathrm{pos}}|
  \leq 2\cdot\codigo{BASE\_ORDER\_SIZE\_XRP}
\end{equation}

Fuera de esos límites la posición se considera no reconciliable y el bot
se detiene y pregunta (§0.2). Además: valida
\codigo{PERCENT\_PRICE} y \codigo{requiredMarginPercent}, consulta el
\emph{mark price}, fija \codigo{set\_leverage(MAX\_LEVERAGE\_USED)} y
exige confirmación interactiva (\codigo{CONFIRMAR}); la bandera
\codigo{--yes} solo se usa bajo supervisión directa o CI. Operar exige
\codigo{config.REAL=True}, \codigo{EXPOSURE\_LEVEL $\geq$ 1} y
\codigo{BUDGETS\_CONFIRMED=True}.

\subsection{Aislamiento en la VPS}
En la misma VPS corren otros bots en producción (grid bot, bot A-S en
BTC/USDC). El proceso nuevo corre aislado: venv propio, logs propios,
sin compartir puertos, archivos ni el límite de rate-limit (§0.3). No
tocar ni reiniciar los procesos de los otros bots.

% =====================================================================
% APÉNDICE A — GLOSARIO
% =====================================================================
\appendix
\section{Glosario}

\begin{description}
  \item[Ask / Bid] Mejor precio de venta / de compra del libro.
  \item[Futures USDS-M] Contratos de futuros liquidados en stablecoin
  (USDC en este caso).
  \item[GTX] \emph{Good 'til cancelled, post-only}: orden maker que se
  cancela si ejecutaría como taker.
  \item[Inventario] Cantidad de XRP acumulada por el bot; positiva =
  larga, negativa = corta.
  \item[Kill switch] Mecanismo de emergencia que cancela todas las
  órdenes y reconcilia la posición ante condiciones anómalas (§13).
  \item[Largo / Corto] Posición de compra / de venta.
  \item[Mark price] Precio de referencia usado para liquidaciones y
  cálculo de nocional.
  \item[Post-only] Ver \emph{GTX}.
  \item[Round trip] Posición pre-existente considerada válida si
  cumple los límites de pre-flight (§7.2).
  \item[Skew] Sesgo aplicado a los precios de cotización según
  inventario (§3.3).
  \item[Testnet] Red de pruebas de Binance con fondos simulados.
  \item[Tick size / Step size] Incremento mínimo de precio / de
  cantidad de un símbolo.
\end{description}

% =====================================================================
% APÉNDICE B — CONSTANTES CLAVE
% =====================================================================
\section{Constantes clave de configuración}

\begin{center}
\begin{longtable}{lll}
\toprule
\textbf{Constante} & \textbf{Valor} & \textbf{Rol} \\
\midrule
\endhead
\codigo{SYMBOL} & \codigo{"XRPUSDC"} & Instrumento (parámetro, no hardcodeado) \\
\codigo{REAL} & \codigo{True} & Mainnet autorizada 10/08/2026 (§0.2) \\
\codigo{EXPOSURE\_LEVEL} & $0$--$3$ & Nivel de exposición (tabla §4.3) \\
\codigo{SAMPLING\_INTERVAL\_SEC} & $5.0$ & Intervalo de muestreo de retornos \\
\codigo{INTERVALS\_PER\_MINUTE} & $12$ & Reescala de volatilidad (§3.1) \\
\codigo{VOLATILITY\_WINDOW\_SEC} & — & Ventana de estimación de vol \\
\codigo{CYCLE\_INTERVAL\_SEC} & — & Duración del ciclo de cotización \\
\codigo{MAX\_DAILY\_LOSS\_USDC} & $10.0$ & Límite diario de pérdida (§0.4) \\
\codigo{MAX\_POSITION\_NOTIONAL\_USDC} & — & Límite de nocional \\
\codigo{MAX\_DRAWDOWN\_PCT} & — & Límite de drawdown \\
\codigo{MAX\_LEVERAGE\_USED} & — & Apalancamiento máximo \\
\codigo{BASE\_ORDER\_SIZE\_XRP} & — & Tamaño base por lado \\
\codigo{MAX\_ORDER\_SIZE\_XRP} & — & Tamaño máximo de orden (clamp §3.4) \\
\codigo{SIMULATION\_QUOTE\_MULTIPLIER} & — & Multiplicador de simulación (Nivel 0) \\
\codigo{FUNDING\_RATE\_PER\_8H} & — & Tasa de funding usada en NetPnL \\
\bottomrule
\end{longtable}
\end{center}

Nota: los valores marcados con ``—'' corresponden a la configuración
vigente en \codigo{strategy/config.py}; consultar ese archivo como
fuente de verdad antes de modificar cualquier límite (§0.4).

\end{document}
